\section{Classical Controller Tuning: Insights and Lessons}
\label{appendix:tuning_insights}

This appendix documents the tuning process for each classical
controller, including approaches that failed and the reasoning
behind final parameter choices.
All tuning uses the full nonlinear dynamics with first-order motor
model ($\tau_v = 0.15$\,s), 1-step command delay, and RK4 integration
at 20~Hz. (S-curve was used in earlier ablations only; the first-order
model matches what is used for final evaluation.)

\subsection{General Methodology}

We tune the classical and predictive controllers with
CMA-ES~\cite{hansen2016cma}, scaling the population size and evaluation
budget to each controller's cost per evaluation: 40 and 1500 evaluations
for PD, 32 and 980 for LQR, 36 and 1200 for SMC, and 20 and 500 for MPC
and NMPC, whose per-evaluation solves are far more expensive. The deployed LQR+WO is an
exception: its gain is the PD-equivalent one (the deployed PD law's gains with
the cart-state terms zeroed) rather than the CMA-ES or Riccati LQR solution,
which fail on hardware (see the LQR subsection below). All share the composite CMA-ES objective $J$ (integrated squared error in
$\theta$ and $u$, a Gaussian-gated ISE in $\dot\theta$, and a
boundary-violation count) defined in Appendix~\ref{appendix:implementation}.
We optimize over
7 stratified initial conditions, including 2 same-side cases with the ball
displaced toward the wall, using 30\,s trials, and evaluate the tuned
controllers on 50 stratified initial conditions with a settling band of
0.01\,rad held for 1.0\,s.

\subsection{PD}

The deployed PD controller uses two gains, $K_p=19.78$ and $K_d=1.37$, together with the
wall override ($d_w=0.12$~m, $u_\text{ov}=0.5$). We set the wall margin at
0.12\,m empirically; it serves as a heuristic braking-distance reference,
close to the velocity-lag product
$V_\text{max} \times \tau_v = 0.9 \times 0.15 = 0.135$~m, which gives a
rough upper bound.

Three additions that seemed reasonable did not survive tuning. A cart-position
centering gain ($u \mathrel{+}= K_x \cdot x$, $K_x < 0$) conflicts with the $K_p \cdot \theta$ term,
since the two oppose in sign when the ball is displaced toward the wall, so we
removed it entirely. Integral action ($K_i > 0$) causes windup during the
wall-override intervals, and CMA-ES consistently converges to
$K_i \approx 0$, which is why we use PD rather than PID. Vanilla PD without
the override reaches 84\% in simulation but only 18\% on hardware
(Table~\ref{tab:failed_variants}): when the ball is off-center it commands the
cart toward the ball, and if the ball responds slowly, through inertia or
friction, the cart crosses the full track and reaches the rail limit before
the ball recenters, a runaway from which PD without the wall override cannot
recover.

\subsection{LQR}

Our LQR computes its gain from the discrete-time algebraic Riccati equation
(DARE): a state weighting $Q$ and control weighting $R$ go in and the optimal
feedback gain comes out. We tried both a hand-set
$Q = \text{diag}([0.001,\, 0.004,\, 18.7,\, 1.76])$, $R = 0.017$ (giving
$K_\theta \approx -20.7$) and a from-scratch CMA-ES search over $Q, R$ against
our usual objective (giving a more aggressive $K_\theta \approx -34$). Both are
optimal-control gains, and both carry cart-position and cart-velocity feedback.

The optimal-control gain did not transfer in our runs. The CMA-ES gain, even
with the wall override, reached only 16\% SR (8/50) and failed from the roomy
center starts too (0 of 10), so what we see looks more like a settling problem
than a boundary one. In the trajectories, the cart-state feedback keeps the
controller active even once the ball is balanced: over the timesteps where the
ball is already inside the $0.01$~rad success band, the command is dominated by
the cart-position term (the cart sits about $0.6$~m off-center at those moments
on average), and the optimal gain saturates the $\pm 0.9$~m/s actuator on more
than half of them, against essentially none for the PD-equivalent gain. Those
cart-centering commands appear to knock the ball back out before it can hold the
band for the required second, the same tension as the Q-center variant below.
The hand-set Riccati gain, run without the override, does no better (3/50).

The gain that did transfer is the PD law used directly as state feedback: we
take the deployed PD gains ($K_\theta = -19.78$, $K_{\dot\theta} = -1.37$) as
$K$, with the cart entries at zero. This gain carries no cart-position or
cart-velocity feedback at all, so on its own it does nothing to keep the cart
off the rail: run without the override it reaches only 20\% SR (10/50), failing
near-uniformly across the difficulty strata as the cart drifts out. Adding the
wall override ($d_w=0.12$~m, $u_\text{ov}=0.5$), with the gain otherwise
unchanged, supplies exactly that missing rail protection and lifts the same
controller to 92\% SR (46/50); this is what the deployed LQR+WO applies.
What we take from this, on this platform rather than as a general claim, is that
reaching a reliable LQR here took two things together: using the PD gains as the
state-feedback gain, so the cart authority the optimal cost assigns does not
fight ball stabilization, and adding the override to handle the rail that this
gain ignores. The 20\%-versus-92\% pair isolates the override's effect with the
gain held fixed; the 16\% optimal gain shows that simply restoring cart
authority is not the fix.

One further variant shows why the cost cannot substitute for the override.
Trying to prevent wall collisions through the cost alone, by raising the
cart-position weight until $K_0 > 2$, makes the centering term fight ball
stabilization and the response oscillates; this forced Q-center variant
reaches only 28\% on hardware. That is a fundamental limitation of a single
linear gain vector, which cannot simultaneously optimize same-side recovery
and large-angle stability.

\paragraph{Override contribution.}
The wall-override, not the underlying control law, does much of the
recovery work for every classical controller on this platform, not just
LQR. Across the 50 trials per controller, the cart reaches the override
trigger zone ($|x| > 0.66$\,m) in 75--98\% of successful trials (85\%
for LQR+WO, 75\% for PD+WO, 98\% for SMC+WO, 80\% for MPC+WO) and, once
there, spends $\sim$31--47\% of its steps inside it. Few trials end at
track center (Fig.~\ref{fig:wall_override_dep}): most successful
trajectories invoke the override and finish away from center, consistent
with boundary handling, rather than linear-model fidelity or
sliding-surface design, doing much of the recovery.

\begin{figure}[H]
  \centering
  \includegraphics[width=0.95\columnwidth]{figures/wall_override_dependence.eps}
  \caption{Wall-override dependence across the four WO classical
    controllers (successful trials only). (a) Override firing rate and
    mean active-step fraction. (b) End-of-trial cart-position
    histogram; few trials end at center, indicating that successful
    trials typically invoke the override rather than settling under the
    nominal control law alone.}
  \label{fig:wall_override_dep}
\end{figure}

\subsection{SMC}

The deployed SMC uses a first-order sliding surface
$\sigma = \dot{\theta} + \lambda\theta$ with a smooth saturation boundary
layer $\phi$ and the same wall override as PD and LQR. Its best variant takes
a relatively wide boundary layer ($\phi=0.806$) to avoid chattering at the
20\,Hz hardware rate, and reaches 92\% on hardware: the override prevents
irreversible wall impacts, while the wide boundary layer keeps the
sensor-noise-induced switching manageable.

The failure modes trace back to the same two pressures. Without boundary
handling, the nominal sliding surface regulates ball angle but does not center
the cart, so residual model-cancellation bias and friction mismatch let the
cart drift toward the rail, leaving the basic variant at 12\% on hardware.
Thin boundary layers make the controller behave more like a discontinuous
sliding controller and chatter severely under noisy ToF feedback, with more
than 500 sign reversals per episode in preliminary tests, so hardware
stability forces the wide-boundary-layer regime. Aggressive cart-centering or
damping gains reduce wall drift but destabilize ball regulation at large
angles, which is why the final SMC+WO variant leans on the override for
boundary recovery rather than trying to solve boundary avoidance inside the
sliding surface. Finally, because SMC's equivalent control depends on accurate
dynamics, the residual sim-to-real gap, consistent with the higher-order motor
response (Appendix~\ref{appendix:simulation_validation}) and sensor staleness
(Appendix~\ref{appendix:sensor_characterization}), is a natural explanation for
why SMC+WO settles more slowly on hardware than in first-order simulation
despite its high success rate.

\subsection{Linear MPC}

Linear MPC uses a short horizon $N=10$ with
$Q = \text{diag}([0.001,\, 0,\, 40,\, 0])$ and $R = 0.017$, together with the
wall override ($d_w=0.12$~m, $u_\text{ov}=0.5$) and the QP cart constraints
\emph{disabled} (the cart-constraint rows are omitted from the QP). Four
things went wrong before we reached this configuration.

Penalizing $\dot{\theta}$ with a positive weight inverts the command. Writing
the linearized discrete-time model as $x_{k+1}=A x_k + B u_k$ with state
$x=[x_\text{cart},\,\dot{x}_\text{cart},\,\theta,\,\dot{\theta}]^\top$ and
input matrix $B=(B_1,B_2,B_3,B_4)^\top$, the $\dot\theta$ input-sensitivity
row $B_4$ is negative near $\theta=0$: a positive cart-velocity command
produces a negative angular acceleration. The $\dot\theta$-penalty is the
fourth diagonal entry $Q_{44}$, and empirically a positive $Q_{44}$ acting
through the negative $B_4$ drives the QP to command the wrong sign, so the
system diverges immediately; setting $Q_{44}=0$ removes the inversion. A long
horizon is the second problem: over many prediction steps the linearization
error accumulates, and at large angles ($\theta > 0.06$)
the linear model underestimates the restoring torque and plans too weak a
correction, so the QP under-responds; $N=30$ passes only 1 of 7 tuning initial conditions.
The QP cart constraint fails on hardware for a related reason: the linear
model underestimates braking distance (first-order lag versus the S-curve jerk
limits), so when the cart is near the rail with positive velocity the QP
becomes infeasible, because no admissible input sequence keeps the predicted
cart positions within the bound, and the fallback $u=0$ lets the ball escape: the natively
constrained variant reaches only 44\% on hardware. Removing the QP cart
constraint and relying on the wall override instead yields 80\%.
Finally, too small an input penalty $R$ (against the tuned $R=0.017$) makes the
controller over-trust the linear model's predicted optimum, which diverges
quickly from reality on hardware and produces oscillatory commands that rarely
settle within the trial.

\subsection{NMPC}

NMPC solves a full nonlinear program in CasADi~\cite{casadi} over the
friction-augmented dynamics, with a short horizon $N=10$,
$Q = \text{diag}([0,\, 0,\, 9,\, 0])$, $R = 0.006$, hard cart constraints in
the NLP, and no wall override; it reaches 90\% on hardware (45 of 50 trials).

Its failures come from optimizing too far ahead of reliable observations. A
long horizon $N=30$ balances 100\% in simulation but passes only 1 of 7
conditions in a preliminary hardware screening, because the NLP optimizes a
longer trajectory from the current, possibly stale estimate and its
predictions degrade once the ball enters the sensor blind zone
($|\theta| > 0.071$~rad) and model error accumulates over the horizon; the
shorter $N=10$ optimizes over fewer steps and is the configuration we deploy.
Both horizons re-solve every 50~ms and apply only the first planned action.
Adding a wall override actually hurts NMPC, because it replaces that first
planned action and invalidates the warm-started trajectory, producing
oscillatory behavior. The first-order-versus-S-curve motor mismatch matters less than it
might: with $N=10$ the solver replans every 50\,ms and applies only the first
planned action, which limits error accumulation, whereas at $N=30$ the prediction quality
degrades and the solver returns unreliable plans.

\subsection{Cross-Controller Insights}

A few lessons cut across all five controllers. The first is that the wall
override is controller-dependent: it helps PD, LQR, SMC, and linear MPC, which
lack reliable boundary handling of their own, but it is counterproductive for
NMPC, whose planned trajectories it disrupts; across the four controllers it
helps, the gain ranges from +36\,pp for MPC to +80\,pp for SMC
(Fig.~\ref{fig:wall_override_ablation}). For MPC this comparison also removes
the QP cart constraint, so its figure reflects the full configuration change
rather than the override alone. The second concerns the sensor
blind zone, which is a ball-side artifact, separate from the cart's rail
limit: when the ball rolls to the ends of the arc ($|\theta| > 0.071$\,rad) it
comes too close to a sensor and the reading saturates. Its contribution to
hardware failure is established directly by the blind-zone ablation, where
modeling it lifts PPO~(BZ+DR) about 8\,pp over its no-blind-zone counterpart
in the development regime (Appendix~\ref{appendix:results}; the matched
first-order-lag checkpoint gives 6\,pp, Appendix~\ref{appendix:rl_fairness}),
rather than inferred from which cart strata fail.
This is a difficulty distinct from the cart running out of rail, even though a
hard recovery can provoke both at once, and the controllers least affected on
hardware are those that recover from the rail quickly (LQR+WO) or that learn
the sensor artifact.

The remaining lessons are about planning. Shorter horizons beat longer ones on
hardware: reducing $N$ from 30 to 10 improves success dramatically for both
MPC and NMPC, because stale observations corrupt long-horizon plans while
reactive replanning tolerates the sensor limitations. More broadly, on this
platform reactivity appears to matter as much as model accuracy. LQR+WO, with a
linear model and no horizon planning, reaches 92\% on hardware, whereas NMPC,
with a full nonlinear model and 10-step planning, reaches 90\% without an
override and is degraded by adding one. No pairwise success-rate difference
among the deployed controllers is statistically significant
(Appendix~\ref{appendix:per_trial}), so we read this as a consistent pattern
rather than a strict ranking: a controller that responds to the \emph{current}
observation is at least competitive with sophisticated planners whose plans are
disrupted by external constraint handling. A final note, from NMPC, is that the
S-curve simulation can be harder than hardware. NMPC $N=10$ scores markedly
lower in S-curve simulation, its internal first-order model running against an
S-curve plant, than the 90\% it reaches on hardware; the S-curve simulation
thus appears to \emph{overestimate} the sim-to-real gap for this short-horizon
predictive controller.

\begin{figure}[H]
  \centering
  \includegraphics[width=0.95\columnwidth]{figures/wall_override_ablation.eps}
  \caption{Effect of wall override on hardware success rate across four
    classical controllers. The improvement ranges from +36\,pp (MPC) to
    +80\,pp (SMC).}
  \label{fig:wall_override_ablation}
\end{figure}
